% !PRINTED 370
\weberpage{370}

Der Unterschied (5) sinkt mit unendlich wachsendem $n$
unter jede Grenze und ist, da $x_n > a_n$, also $Q_n x_n + Q_{n-1} > Q_{n+1}$
ist, dem absoluten Werth nach kleiner als
\weq{6}{\frac{1}{Q_n Q_{n+1}}}

und um so mehr kleiner als $1 : Q_n^2$.

Die Zahlen der Reihe (3) nähern sich also abnehmend, die
der Reihe (4) zunehmend der Grenze $x$. Der Ausdruck (6) giebt
ein Maass für den Fehler, den man begeht, wenn man $x$ durch
den Näherungsbruch $P_n : Q_n$ ersetzt.

Die Näherungsbrüche sind also angenäherte
Ausdrücke von Irrationalzahlen durch rationale
Brüche.

Dass diese Näherungsbrüche bei gleichem Grade der Annäherung
 an die Irrationalzahl $x$ die möglichst einfachen sind,
das wird durch folgenden Satz ausgedrückt:

\gesperrt{IV. Es lässt sich zwischen den zwei rationalen
Brüchen}
\[
\frac{P_n}{Q_n}, \quad \frac{P_{n-1}}{Q_{n-1}}
\]
\gesperrt{kein anderer rationaler Bruch einschieben,
dessen Nenner kleiner als $Q_n$ oder auch nur
gleich $Q_n$ ist.}

Angenommen, es liege der rationale Bruch $M : N$ zwischen
den Näherungsbrüchen $P_n : Q_n$ und $P_{n-1} : Q_{n-1}$. Dann ist die
Differenz
\[
\frac{P_n}{Q_n} - \frac{P_{n-1}}{Q_{n-1}}
\]
absolut grösser und vom selben Vorzeichen wie
\[
\frac{M}{N} - \frac{P_{n-1}}{Q_{n-1}},
\]
also
\[
(-1)^n \left(\frac{P_n}{Q_n} - \frac{P_{n-1}}{Q_{n-1}}\right) > (-1)^n \left(\frac{M}{N} - \frac{P_{n-1}}{Q_{n-1}}\right).
\]

Multiplicirt man beiderseits mit $N Q_{n-1}$, so folgt nach (1)
\[
\frac{N}{Q_n} > (-1)^n (M Q_{n-1} - N P_{n-1}),
\]
und da rechts eine positive ganze Zahl steht, die also mindestens
gleich 1 ist, so folgt
\[
N > Q_n.
\]
